%% TODO: Add a small paragraph to tell what this chapter is about
This chapter delves into the existing literature and technological advancements that form the foundation of the UrbanPulse platform. It evaluates current methodologies in real-time urban monitoring, analyzes the architectural shift toward Data Lakehouses, and identifies the research gaps that necessitate a dual-signal anomaly detection framework.
%% TODO: Make sure to use \textbf or \textit for highlighting keywords, and \cite{} to cite the corresponding quotations
\section{Current Advancements}
The landscape of Intelligent Transportation Systems (ITS) has evolved significantly with the integration of high-throughput messaging systems. Recent studies emphasize the use of \textbf{Apache Kafka} and \textbf{Redpanda} for low-latency ingestion of heterogeneous urban signals, such as traffic flow and weather telemetry~\cite{urban-pulse-platform}. In the domain of data storage, the industry is transitioning from traditional data warehouses to the \textbf{Medallion Lakehouse architecture}. Utilizing table formats like \textbf{Apache Iceberg} allows for transactional integrity and schema evolution, which are critical for maintaining long-term urban datasets. Furthermore, state-of-the-art anomaly detection now leverages a combination of \textbf{statistical baselines} and \textbf{unsupervised learning}, such as Isolation Forest, to handle the high dimensionality of city-scale data.
\section{Research Gap}
Despite these advancements, a significant \textit{semantic gap} persists between raw telemetry ingestion and actionable intelligence. Most existing systems process traffic data in isolated batches, leading to high latency in incident response. There is a lack of integrated frameworks that ensure \textbf{stream-to-batch consistency}, where real-time anomaly signals are validated against historical neighborhood rhythms stored in a Lakehouse~\cite{5steps-gaps}. Furthermore, many platforms struggle with the "cold-start" problem in low-variance traffic conditions, often resulting in high False Positive Rates (FPR). UrbanPulse addresses these gaps by implementing a unified \textbf{dual-signal detector} that reconciles instant Z-score deviations with historical machine learning inference.

%% TODO: Make sure to use \subsection{} and \subsubsection{} for smaller sections inside a larger sections
\section{Theoretical Background}
\subsection{The Medallion Lakehouse Architecture}
The core data strategy of this research follows the Medallion pattern, organized into three distinct layers: \textbf{Bronze} (raw ingestion), \textbf{Silver} (filtered and cleaned), and \textbf{Gold} (aggregated business-level features). This structure, supported by \textbf{Project Nessie} for catalog versioning, ensures that data remains reproducible and queryable at any point in time as illustrated in the system's data flow.

\begin{figure}[ht]
	\centering
	% Đổi tên file thành hình kiến trúc thực tế của bạn nếu có
	\includegraphics[width=1\linewidth]{images/system-arch.png} 
	\caption{Conceptual overview of the Medallion Lakehouse architecture in UrbanPulse.}
	\label{fig:lakehouse}
\end{figure}

\subsubsection{Streaming-First Processing}
Unlike traditional ETL, a \textbf{streaming-first approach} minimizes the temporal gap between event occurrence and dashboard visualization. By using micro-batching in Python and pushing updates via \textbf{Server-Sent Events (SSE)}, the system achieves near-real-time updates for end-users. A comparison of the architectural components used in UrbanPulse versus traditional systems is provided in Table~\ref{tab:1}.

%% TODO: Adjust the size of the figure by using [widht=.x\linewidth] (x is a fraction) to fit within the page width. Then rename to your picture file name, add the caption and a label
\begin{table}[ht]
	\centering
	\caption{Comparison of different architectural methods (\protect\cmark: YES, \protect\xmark: NO).}
	\label{tab:1}
	\resizebox{\textwidth}{!}{%
		\begin{tabular}{lcccc}
			\hline
			\textbf{Feature}          & \textbf{UrbanPulse} & Traditional Batch & Lambda Architecture & Lambda-Plus \\ \hline
			Real-time Processing      & \cmark              & \xmark            & \cmark              & \cmark      \\ 
			Schema Evolution (Iceberg)& \cmark              & \xmark            & \xmark              & \cmark      \\ 
			Dual-Signal Detection     & \cmark              & \xmark            & \xmark              & \xmark      \\ 
			Time-Travel Queries       & \cmark              & \cmark            & \xmark              & \cmark      \\ \hline
		\end{tabular}%
	}
\end{table}

\subsection{Anomaly Detection Logic}
UrbanPulse utilizes a dual-signal framework to identify traffic incidents. The statistical signal is calculated using the Z-score formula to detect immediate spikes in congestion ratio $r$ relative to the moving average $\mu$ and standard deviation $\sigma$:

\begin{equation}
    Z = \frac{r - \mu}{\sigma}
\end{equation}

Simultaneously, an \textbf{Isolation Forest} model operates on cyclical temporal features (e.g., hour of day, day of week) to identify non-recurrent patterns that simple statistics might miss. The relationship between the input features and the anomaly score is modeled through a multi-dimensional mapping:

\begin{align}
	f : \mathbb{R}^n \times \mathcal{T} & \rightarrow [-1, 1] \\
	(\mathbf{x}, t)                     & \mapsto s          
\end{align}

where $\mathbf{x}$ represents traffic metrics, $t$ is the cyclical time component, and $s$ is the resulting anomaly score.
